% ============================================================
% INTRODUCTION GÉNÉRALE
% ============================================================
\chapter*{\centering Introduction générale}
\addcontentsline{toc}{chapter}{Introduction générale}
\markboth{Introduction générale}{Introduction générale}

La transformation numérique du secteur traiteur impose aujourd'hui une meilleure maîtrise des flux d'information, de la relation client et des processus opérationnels. Une gestion reposant uniquement sur des pratiques manuelles limite la réactivité, augmente le risque d'erreurs et complique le suivi des commandes et des événements.

Au sein de \textbf{Assiette Gala Sfaxienne}, l'analyse de l'existant a mis en évidence plusieurs difficultés: dispersion des données, traitement peu fluide des demandes de devis, suivi partiel des commandes et absence d'outil d'assistance client automatisé. Dans ce contexte, ce projet de fin d'études vise la conception et la réalisation d'une plateforme web intégrée, adaptée aux besoins réels de l'entreprise.

La solution proposée repose sur une architecture \textbf{microservices} articulée autour d'une passerelle API. Les domaines fonctionnels principaux (authentification, catalogue, commandes, événements et intelligence artificielle) sont isolés en services spécialisés afin d'améliorer la maintenabilité, la scalabilité et l'évolutivité du système.

Le projet intègre également un module d'intelligence artificielle fondé sur une approche \textbf{RAG} (Retrieval-Augmented Generation) pour assister les utilisateurs et fournir des réponses contextualisées à partir des données métier. Ce choix permet d'améliorer l'expérience utilisateur tout en conservant la cohérence des informations fournies. Le stockage vectoriel repose sur \textbf{FAISS} (Facebook AI Similarity Search), garantissant une recherche de similarité performante et un déploiement simplifié sur Windows.

La réalisation a été conduite selon la méthodologie \textbf{Scrum}, avec une planification en cinq sprints et une validation progressive des fonctionnalités.

Le présent manuscrit est organisé comme suit:

\begin{itemize}
    \item Le \textbf{premier chapitre} (Étude Préalable) ouvre ce recueil en présentant l'organisme d'accueil, l'examen des logiciels du marché, ainsi que le cadre méthodologique Scrum et la planification.
    \item Le \textbf{deuxième chapitre} (Architecture) détaille les fondations techniques du projet, l'approche microservices et les choix technologiques retenus.
    \item Le \textbf{troisième chapitre} décrit le Sprint~1, centré sur la gestion des utilisateurs et la sécurisation de l'authentification.
    \item Le \textbf{quatrième chapitre} s'attarde sur le Sprint~2, consacré à la modélisation du catalogue culinaire et à la gestion de la logistique.
    \item Le \textbf{cinquième chapitre} détaille le Sprint~3, mettant en lumière l'expérience client à travers la gestion des événements sur mesure et le circuit de commande.
    \item Le \textbf{sixième chapitre} résume le Sprint~4, dédié au pilotage de l'entreprise via le Back-Office, l'analyse des KPI et la génération de documents comptables.
    \item Le \textbf{septième chapitre} expose le Sprint~5, qui intègre la couche d'Intelligence Artificielle (chatbot RAG, prévisions de stock, optimisation en cuisine et reconnaissance OCR).
\end{itemize}

Une conclusion générale clôture ce rapport en synthétisant les résultats obtenus, les limites identifiées et les perspectives d'amélioration.

